\documentclass[12pt]{article}
\usepackage{amsmath,amssymb,amsfonts}
\usepackage[margin=1in]{geometry}
\begin{document}
\begin{center}
{\large\bfseries 22nd Irish Mathematical Olympiad\\9 May 2009}
\end{center}
\bigskip\noindent\textbf{Paper 1}\bigskip
\begin{enumerate}
\item Hamilton Avenue has eight houses. On one side of the street are the houses numbered $1, 3, 5, 7$ and directly opposite are houses $2, 4, 6, 8$ respectively. An eccentric postman starts deliveries at house $1$ and delivers letters to each of the houses, finally returning to house $1$ for a cup of tea. Throughout the entire journey he must observe the following rules. The numbers of the houses delivered to must follow an odd-even-odd-even pattern throughout, each house except house $1$ is visited exactly once (house $1$ is visited twice) and the postman at no time is allowed to cross the road to the house directly opposite. How many different delivery sequences are possible?

\item Let $ABCD$ be a square. The line segment $AB$ is divided internally at $H$ so that $|AB|\cdot|BH| = |AH|^2$. Let $E$ be the mid point of $AD$ and $X$ be the midpoint of $AH$. Let $Y$ be a point on $EB$ such that $XY$ is perpendicular to $BE$. Prove that $|XY| = |XH|$.

\item Find all positive integers $n$ for which $n^8 + n + 1$ is a prime number.

\item Given an $n$-tuple of numbers $(x_1, x_2, \ldots, x_n)$ where each $x_i = +1$ or $-1$, form a new $n$-tuple
\[
(x_1 x_2,\; x_2 x_3,\; x_3 x_4,\; \ldots,\; x_n x_1),
\]
and continue to repeat this operation. Show that if $n = 2^k$ for some integer $k \ge 1$, then after a certain number of repetitions of the operation, we obtain the $n$-tuple
\[
(1, 1, \ldots, 1).
\]

\item Suppose $a$, $b$, $c$ are real numbers such that $a + b + c = 0$ and $a^2 + b^2 + c^2 = 1$. Prove that
\[
a^2 b^2 c^2 \;\le\; \frac{1}{54},
\]
and determine the cases of equality.
\end{enumerate}

\newpage
\begin{center}
{\large\bfseries 22nd Irish Mathematical Olympiad\\9 May 2009}
\end{center}
\bigskip\noindent\textbf{Paper 2}\bigskip
\begin{enumerate}
\setcounter{enumi}{5}
\item Let $p(x)$ be a polynomial with rational coefficients. Prove that there exists a positive integer $n$ such that the polynomial $q(x)$ defined by
\[
q(x) = p(x+n) - p(x)
\]
has integer coefficients.

\item For any positive integer $n$ define
\[
E(n) = n(n+1)(2n+1)(3n+1)\cdots(10n+1).
\]
Find the greatest common divisor of $E(1)$, $E(2)$, $E(3)$, \ldots, $E(2009)$.

\item Find all pairs $(a,b)$ of positive integers, such that $(ab)^2 - 4(a+b)$ is the square of an integer.

\item At a strange party, each person knew exactly $22$ others.

For any pair of people X and Y who knew one another, there was no other person at the party that they both knew.

For any pair of people X and Y who did not know one another, there were exactly $6$ other people that they both knew.

How many people were at the party?

\item In the triangle $ABC$ we have $|AB| < |AC|$. The bisectors of the angles at $B$ and $C$ meet $AC$ and $AB$ at $D$ and $E$ respectively. $BD$ and $CE$ intersect at the incentre $I$ of $\triangle ABC$.

Prove that $\angle BAC = 60^\circ$ if and only if $|IE| = |ID|$.
\end{enumerate}
\end{document}
